% ============================================================
%  customised-commands.tex  —  merged from all venue variants
%  Variants: (base/generic), ACL, ACM (acmart), LNCS (llncs), Elsevier (elsarticle)
%
%  Auto-detection uses \@ifclassloaded{<class>}{<yes>}{<no>}.
%  All venue-specific logic is resolved at compile time — no
%  manual editing needed when switching document class.
% ============================================================

\makeatletter

% ============================================================
%  CONVENIENCE BOOLEANS  (set once, used throughout)
% ============================================================
\newif\if@isACM    \@isACMfalse
\newif\if@isLNCS   \@isLNCSfalse
\newif\if@isELS    \@isELSfalse
\newif\if@isACL    \@isACLfalse

\@ifclassloaded{acmart}    {\@isACMtrue}{}
\@ifclassloaded{llncs}     {\@isLNCStrue}{}
\@ifclassloaded{elsarticle}{\@isELStrue}{}
% ACL uses several class names across years
\@ifclassloaded{acl}       {\@isACLtrue}{}
\@ifclassloaded{acl2023}   {\@isACLtrue}{}
\@ifclassloaded{acl_natbib}{\@isACLtrue}{}

\makeatother

% -------  optional / review package  -----------------------
% \usepackage{easyReview}

% -------  language  ----------------------------------------
\usepackage[english]{babel}

% -------  geometry  ----------------------------------------
% ACL loads geometry explicitly; other classes manage their own margins.
\makeatletter
\if@isACL
  \usepackage{geometry}
  % \usepackage[letterpaper,top=2cm,bottom=2cm,left=3cm,right=3cm,marginparwidth=1.75cm]{geometry}
\fi
\makeatother

% -------  bibliography  ------------------------------------
% ACM (acmart) has its own citation machinery — skip natbib there.
\makeatletter
\if@isACM\else
  \usepackage{natbib}  % \bibliographystyle{plainnat}
  \let\oldcite\cite
  \renewcommand{\cite}[1]{\citep{#1}}
\fi
\makeatother
% \usepackage[backend=biber,style=numeric-comp,sorting=none]{biblatex}

% -------  graphics & layout  -------------------------------
\usepackage{graphicx}
\usepackage{microtype}
% \usepackage{subfigure}
\usepackage{booktabs}
\usepackage{float}
\usepackage{wrapfig}
\usepackage[normalem]{ulem}

% -------  math  --------------------------------------------
\usepackage{bm}
\usepackage{textcomp}
% amssymb: ACM commented it out (acmart already pulls it in via newtxmath).
\makeatletter
\if@isACM\else
  \usepackage{amssymb}
\fi
\makeatother
\usepackage{amsmath}
% amsthm: llncs defines its own theorem environments — skip it there.
\makeatletter
\if@isLNCS\else
  \usepackage{amsthm}
\fi
\makeatother
\usepackage{amsfonts}
\usepackage{mathtools}
\usepackage{pifont}

% -------  algorithms  --------------------------------------
\usepackage{algorithm}
% ACM uses algpseudocode; all other venues use algorithmic.
\makeatletter
\if@isACM
  \usepackage{algpseudocode}
\else
  \usepackage{algorithmic}
\fi
\makeatother

% LNCS-only: inline comment command for algorithm floats.
\makeatletter
\if@isLNCS
  \newcommand{\algcomment}[2][gray]{%
    \unskip\hfill{\footnotesize\textcolor{#1}{// #2}}%
  }
\fi
\makeatother

% -------  fonts & symbols  ---------------------------------
% ACM (acmart) manages its own fonts — \usepackage{times} conflicts.
\makeatletter
\if@isACM\else
  \usepackage{times}
\fi
\makeatother
% \usepackage{latexsym}
\usepackage{gensymb}

% -------  table & figure utilities  ------------------------
\usepackage{multirow}
\usepackage{lscape}
\usepackage{subcaption}

% -------  frames, boxes, hyperlinks  -----------------------
\usepackage{mdframed}
\usepackage{hyperref}
\usepackage{textgreek}

% -------  color & table color  -----------------------------
\usepackage{tabularx, colortbl}
\usepackage[dvipsnames]{xcolor}

% ACM-only: TikZ physics shorthand color.
% xcol (steel blue) -> RoyalBlue (dvipsnames)
\makeatletter
\if@isACM
  \colorlet{xcol}{RoyalBlue}
\fi
\makeatother

% ACL table header/row shading — available color names for use in tables:
%   headerpink  -> Salmon!30    headergreen -> LimeGreen!30
%   roworange   -> Peach!50     rowgray     -> lightgray!50
%   textorange  -> BurntOrange

% -------  column types & highlight commands  ---------------
% Color mapping from originals to xcolor dvipsnames:
%   p13 (soft lavender) -> Lavender
%   mygray (near-white) -> lightgray
%   pink                -> Pink
\newcolumntype{a}{>{\columncolor{Lavender}}l}
% \newcolumntype{b}{>{\columncolor{LimeGreen}}c}
\DeclareRobustCommand{\hlpurple}[1]{{\sethlcolor{Lavender}\hl{#1}}}
\DeclareRobustCommand{\hlpink}[1]{{\sethlcolor{Pink}\hl{#1}}}
\DeclareRobustCommand{\hlgray}[1]{{\sethlcolor{lightgray}\hl{#1}}}

% \renewcommand{\alertColor}{\textcolor{Red}}
% \renewcommand{\addColor}{\textcolor{Cyan}}

% -------  cleveref  ----------------------------------------
\usepackage[capitalize,noabbrev]{cleveref}
\crefname{ineq}{Inequality}{Inequalities}
\creflabelformat{ineq}{#2{\upshape(#1)}#3}
% obj / exp cref names used in Elsevier and LNCS variants.
\makeatletter
\if@isELS
  \crefname{obj}{Objective}{Objectives}
  \creflabelformat{obj}{#2{\upshape(#1)}#3}
  \crefname{exp}{Expression}{Expressions}
  \creflabelformat{exp}{#2{\upshape(#1)}#3}
\fi
\if@isLNCS
  \crefname{obj}{Objective}{Objectives}
  \creflabelformat{obj}{#2{\upshape(#1)}#3}
  \crefname{exp}{Expression}{Expressions}
  \creflabelformat{exp}{#2{\upshape(#1)}#3}
\fi
\makeatother

% -------  theorem environments  ----------------------------
% The reptheorem helper lets you restate theorems by label.
% Only used in base/ACL variants.
\makeatletter
\if@isACM\else\if@isELS\else\if@isLNCS\else
  \newtheorem*{rep@theorem}{\rep@title}
  \newcommand{\newreptheorem}[2]{%
  \newenvironment{rep#1}[1]{%
  \def\rep@title{#2 \ref{##1}}%
  \begin{rep@theorem}}%
  {\end{rep@theorem}}}
\fi\fi\fi
\makeatother

% LNCS defines all theorem environments itself — skip them there.
\makeatletter
\if@isLNCS\else

  \theoremstyle{remark}
  \newtheorem{example}{Example}[]
  \if@isACM\else\if@isELS\else
    \newreptheorem{example}{Example}
  \fi\fi

  \theoremstyle{plain}
  \newtheorem{theorem}{Theorem}[section]
  \if@isACM\else\if@isELS\else
    \newreptheorem{theorem}{Theorem}
  \fi\fi

  \newtheorem{proposition}[theorem]{Proposition}
  \newtheorem{lemma}[theorem]{Lemma}
  \if@isACM\else\if@isELS\else
    \newreptheorem{lemma}{Lemma}
  \fi\fi

  \newtheorem{corollary}[theorem]{Corollary}
  \if@isACM\else\if@isELS\else
    \newreptheorem{corollary}{Corollary}
  \fi\fi

  \theoremstyle{definition}
  \newtheorem{definition}[theorem]{Definition}
  \newtheorem{assumption}[theorem]{Assumption}
  \theoremstyle{remark}
  \newtheorem{remark}[theorem]{Remark}

\fi % end \if@isLNCS\else
\makeatother

% ACM-only theorem types.
\makeatletter
\if@isACM
  \newtheorem{requirement}{Requirement}[]
  \newtheoremstyle{researchquestionstyle}{}{}{}{}{{\bf}}{}{.5em}{}
  \theoremstyle{researchquestionstyle}
  \newtheorem{rquestion}{Research Question:}[]
\fi
\makeatother

% -------  math operators & notation  -----------------------
\usepackage{braket}
\newcommand{\norm}[1]{\left\lVert#1\right\rVert}
\newcommand{\abs}[1]{\left\lvert#1\right\rvert}
\newcommand{\concat}{\mathbin{{+}\mspace{-8mu}{+}}}
\newcommand{\inner}[1]{\left\langle#1\right\rangle}
\usepackage{stmaryrd}
\newcommand{\iverson}[1]{\left\llbracket#1\right\rrbracket}

\DeclareMathOperator*{\argmax}{arg\,max}
\DeclareMathOperator*{\argmin}{arg\,min}

% -------  physics (ACM only)  ------------------------------
\makeatletter
\if@isACM
  \usepackage{physics}
\fi
\makeatother

% -------  TikZ & PGFPlots  ---------------------------------
\usepackage{tikz}
\usetikzlibrary{trees}
\usepackage{tikz-dependency}
\usepackage{tikz-3dplot}
\usetikzlibrary{chains,scopes}
\usetikzlibrary{arrows.meta,automata,positioning}
\usetikzlibrary{shapes.geometric, arrows}
\usetikzlibrary{shapes.geometric, arrows.meta, positioning, calc, fit, shadows}

% ACM-only: enforce LaTeX-style arrow heads globally.
\makeatletter
\if@isACM
  \tikzset{>=latex}
\fi
\makeatother

\usepackage{pgfplots}
% compat=newest and fillbetween used in base/ACL; set conservatively for others.
\makeatletter
\if@isACM\else\if@isLNCS\else\if@isELS\else
  \pgfplotsset{compat=newest}
  \usepgfplotslibrary{fillbetween}
\fi\fi\fi
\makeatother

\pgfplotsset{
  every axis/.append style = {thick},
  tick style = {thick,black},
  %
  % #1 = x, y, or z  |  #2 = the shift value
  /tikz/normal shift/.code 2 args = {%
    \pgftransformshift{%
        \pgfpointscale{#2}{\pgfplotspointouternormalvectorofticklabelaxis{#1}}%
    }%
  },%
  shift/.style = {
    tick align        = outside,
    scaled ticks      = false,
    enlargelimits     = false,
    ticklabel shift   = {#1},
    axis lines*       = left,
    xtick style       = {normal shift={x}{#1}},
    ytick style       = {normal shift={y}{#1}},
    x axis line style = {normal shift={x}{#1}},
    y axis line style = {normal shift={y}{#1}},
  },
  shift/.default = 10pt,
  shift3d/.style = {
    shift=#1,
    ztick style       = {normal shift={z}{#1}},
    z axis line style = {normal shift={z}{#1}},
  },
  shift3d/.default = 10pt,
}

% TikZ flow-chart styles — Elsevier and LNCS variants only.
\makeatletter
\if@isELS
  \tikzstyle{startstop} = [rectangle, rounded corners, minimum width=1cm, minimum height=0.5cm, text centered, draw=black]
  \tikzstyle{io}        = [trapezium, trapezium left angle=70, trapezium right angle=110, minimum width=2cm, minimum height=0.63cm, text centered, draw=black]
  \tikzstyle{process}   = [rectangle, minimum width=2cm,   minimum height=0.5cm, text centered, draw=black]
  \tikzstyle{process_n} = [rectangle, minimum width=1.5cm, minimum height=0.5cm, text centered, text width=1.5cm, draw=black]
  \tikzstyle{decision}  = [diamond, minimum width=1.2cm, minimum height=0.5cm, text width=1.5cm, text centered, aspect=1.5, draw=black]
  \tikzstyle{arrow}     = [->,>=stealth]
\fi
\if@isLNCS
  \tikzstyle{startstop} = [rectangle, rounded corners, minimum width=1cm, minimum height=0.5cm, text centered, draw=black]
  \tikzstyle{io}        = [trapezium, trapezium left angle=70, trapezium right angle=110, minimum width=2cm, minimum height=0.63cm, text centered, draw=black]
  \tikzstyle{process}   = [rectangle, minimum width=2cm,   minimum height=0.5cm, text centered, draw=black]
  \tikzstyle{process_n} = [rectangle, minimum width=1.5cm, minimum height=0.5cm, text centered, text width=1.5cm, draw=black]
  \tikzstyle{decision}  = [diamond, minimum width=1.2cm, minimum height=0.5cm, text width=1.5cm, text centered, aspect=1.5, draw=black]
  \tikzstyle{arrow}     = [->,>=stealth]
  % Anchor-highlight helper (LNCS only)
  \newcommand{\markline}[1]{%
    \tikz[remember picture,baseline] \node (#1) {};%
  }
\fi
\makeatother

% -------  font icons (ACM only)  ---------------------------
\makeatletter
\if@isACM
  \usepackage{fontawesome5}
  % \usepackage{emoji}   % requires LuaLaTeX/XeLaTeX
\fi
\makeatother

% -------  code listings  -----------------------------------
\usepackage{listings}
\lstset{
  language=Python,
%   numbers=left,
%   numberstyle=\tiny,
  stepnumber=1,
  numbersep=5pt,
  tabsize=4,
  basicstyle=\scriptsize\ttfamily,
  columns=fullflexible,
  keepspaces,
%   xleftmargin=\parindent,
  showstringspaces=false
}

% JSON language definition (ACM only).
% Color mapping: numb -> Maroon | punct -> BrickRed | delim -> OliveGreen
\makeatletter
\if@isACM
  \lstdefinelanguage{json}{
    basicstyle=\scriptsize\ttfamily,
    numbers=none,
    stepnumber=1,
    numbersep=5pt,
    showstringspaces=false,
    breaklines=true,
    frame=single,
    backgroundcolor=\color{lightgray!30},
    literate=
     *{0}{{{ \color{Maroon}0}}}{1}
      {1}{{{ \color{Maroon}1}}}{1}
      {2}{{{ \color{Maroon}2}}}{1}
      {3}{{{ \color{Maroon}3}}}{1}
      {4}{{{ \color{Maroon}4}}}{1}
      {5}{{{ \color{Maroon}5}}}{1}
      {6}{{{ \color{Maroon}6}}}{1}
      {7}{{{ \color{Maroon}7}}}{1}
      {8}{{{ \color{Maroon}8}}}{1}
      {9}{{{ \color{Maroon}9}}}{1}
      {:}{{{ \color{BrickRed}{:}}}}{1}
      {,}{{{ \color{BrickRed}{,}}}}{1}
      {\{}{{{ \color{OliveGreen}{\{}}}}{1}
      {\}}{{{ \color{OliveGreen}{\}}}}}{1}
      {[}{{{ \color{OliveGreen}{[}}}}{1}
      {]}{{{ \color{OliveGreen}{]}}}}{1},
  }
  % \lstset{language=json}
\fi
\makeatother

% -------  array & table helpers  ---------------------------
\usepackage{array}
\newcolumntype{H}{>{\setbox0=\hbox\bgroup}c<{\egroup}@{}}
\newcommand{\RomanNumeralCaps}[1]{\MakeUppercase{\romannumeral #1}}
\newcommand{\rnumcap}[1]{\MakeUppercase{\romannumeral #1}}
\newcommand\T{\rule{0pt}{3.4mm}}
\newcommand\B{\rule[-1.5mm]{0pt}{0pt}}

% -------  tcolorbox  ---------------------------------------
\usepackage[most]{tcolorbox}
\newtcolorbox[auto counter]{summary}[1][]{title={\bfseries Summary~\thetcbcounter},enhanced,
  coltitle=black,
  colback=white,
  top=0.3in,
  attach boxed title to top left=
  {xshift=1.5em,yshift=-\tcboxedtitleheight/2},boxrule=0.5pt, sharp corners,
  fonttitle=\bfseries,
  boxed title style={size=small,colback=white,colframe=white},#1}
